\documentclass[aspectratio=169, 10pt, xcolor=table]{beamer}

% --- Configuración del Tema ---
\usetheme{Madrid}
\usecolortheme{dolphin}

% --- Paquetes necesarios ---
\usepackage[utf8]{inputenc}
\usepackage[spanish]{babel}
\usepackage{amsmath, amssymb}
\usepackage{booktabs}
\usepackage{graphicx}
\usepackage{listings}
\usepackage{tikz}
\usetikzlibrary{arrows.meta, positioning, shapes.geometric}
\usepackage{etoolbox}

% --- Estilo para codigo Python ---
\definecolor{codegreen}{rgb}{0,0.6,0}
\definecolor{codegray}{rgb}{0.5,0.5,0.5}
\definecolor{codepurple}{rgb}{0.58,0,0.82}
\definecolor{backcolour}{rgb}{0.95,0.95,0.92}

\lstset{
    language=Python,
    backgroundcolor=\color{backcolour},
    commentstyle=\color{codegreen},
    keywordstyle=\color{blue},
    numberstyle=\tiny\color{codegray},
    stringstyle=\color{codepurple},
    basicstyle=\ttfamily\scriptsize,
    breaklines=true,
    frame=single,
    showstringspaces=false,
    literate={á}{{\'a}}1 {é}{{\'e}}1 {í}{{\'i}}1 {ó}{{\'o}}1 {ú}{{\'u}}1 {ñ}{{\~n}}1 {ü}{{\"u}}1
}

% --- Pie de pagina personalizado ---
\makeatletter
\setbeamertemplate{footline}{
  \leavevmode%
  \hbox{%
  \begin{beamercolorbox}[wd=.333333\paperwidth,ht=2.25ex,dp=1ex,center]{author in head/foot}%
    \usebeamerfont{author in head/foot}\insertshortauthor
  \end{beamercolorbox}%
  \begin{beamercolorbox}[wd=.333333\paperwidth,ht=2.25ex,dp=1ex,center]{title in head/foot}%
    \usebeamerfont{title in head/foot}Sesion 3
  \end{beamercolorbox}%
  \begin{beamercolorbox}[wd=.333333\paperwidth,ht=2.25ex,dp=1ex,right]{date in head/foot}%
    \usebeamerfont{date in head/foot}CALCULO y ALGEBRA LINEAL \hfill \insertframenumber/\inserttotalframenumber \hspace*{0.3cm}
  \end{beamercolorbox}}%
  \vskip0pt%
}
\makeatother

% --- Datos de la portada ---
\title[SESION 3 CALCULO]{\textbf{Sesion 3}}
\subtitle{Operaciones con Tensores}
\author[Prof. C. Sanchez]{Carlos Cesar Sanchez Coronel}
\institute{\textbf{CALCULO Y ALGEBRA LINEAL PARA IA}}
\date{}

\setbeamertemplate{navigation symbols}{}

\AtBeginSection[]{
  \begin{frame}
    \frametitle{Contenido}
    \tableofcontents[currentsection]
  \end{frame}
}

\begin{document}

\begin{frame}
    \titlepage
\end{frame}

\begin{frame}{Contenido}
    \tableofcontents
\end{frame}

\section{Operaciones Elemento a Elemento}

\begin{frame}{Suma, resta y producto Hadamard}
    \begin{columns}[t]
        \column{0.5\textwidth}
        \textbf{Suma:}
        \[ (A+B)_{ij} = A_{ij} + B_{ij} \]
        \textbf{Resta:}
        \[ (A-B)_{ij} = A_{ij} - B_{ij} \]
        \column{0.5\textwidth}
        \textbf{Producto Hadamard:}
        \[ (A \odot B)_{ij} = A_{ij} \cdot B_{ij} \]
    \end{columns}
    \vspace{0.3cm}
    \begin{exampleblock}{Ejemplo con matrices}
        \[ A = \begin{pmatrix}1&2\\3&4\end{pmatrix},\quad B = \begin{pmatrix}5&6\\7&8\end{pmatrix} \]
        \[ A+B = \begin{pmatrix}6&8\\10&12\end{pmatrix},\quad A\odot B = \begin{pmatrix}5&12\\21&32\end{pmatrix} \]
    \end{exampleblock}
\end{frame}

\begin{frame}[fragile]{En Python con NumPy}
\begin{lstlisting}
import numpy as np
A = np.array([[1,2],[3,4]])
B = np.array([[5,6],[7,8]])
print(A + B)
print(A - B)
print(A * B)
\end{lstlisting}
\end{frame}

\section{Broadcasting}

\begin{frame}{Que es broadcasting?}
    \begin{itemize}
        \item Permite operar entre tensores de diferentes formas, estirando el mas pequeño sin copiar datos.
        \item \textbf{Reglas} (desde el ultimo eje hacia el primero):
        \begin{enumerate}
            \item Las dimensiones se comparan de atras hacia adelante.
            \item Dos dimensiones son compatibles si son iguales o una de ellas es 1.
            \item Si una dimension falta en un tensor, se trata como 1.
        \end{enumerate}
    \end{itemize}
\end{frame}

\begin{frame}{Ejemplos de broadcasting}
    \begin{itemize}
        \item Escalar + vector: $5 + [1,2,3] = [6,7,8]$.
        \item Vector + matriz: $[1,2,3] + \begin{pmatrix}10&20&30\\40&50&60\end{pmatrix} = \begin{pmatrix}11&22&33\\41&52&63\end{pmatrix}$.
        \item Vector columna + matriz: $\begin{pmatrix}1\\2\end{pmatrix} + \begin{pmatrix}10&20\\30&40\end{pmatrix} = \begin{pmatrix}11&21\\32&42\end{pmatrix}$.
    \end{itemize}
\end{frame}

\begin{frame}[fragile]{Broadcasting en NumPy}
\begin{lstlisting}
import numpy as np
a = 5
v = np.array([1,2,3])
print(a + v)

v = np.array([1,2,3])
M = np.array([[10,20,30],[40,50,60]])
print(M + v)

v_col = np.array([[1],[2]])
M = np.array([[10,20],[30,40]])
print(M + v_col)
\end{lstlisting}
\end{frame}

\begin{frame}[fragile]{Aplicacion: Batch Normalization}
    En Batch Normalization, se normaliza cada canal de un lote de imagenes. La media por canal tiene forma $(1,1,1,C)$ y se resta del tensor $(B,H,W,C)$ usando broadcasting.
    
    \begin{exampleblock}{Simulacion simplificada}
\begin{lstlisting}
import numpy as np
batch = np.random.rand(32,32,32,3)
media_por_canal = batch.mean(axis=(1,2), keepdims=True)
batch_norm = batch - media_por_canal
\end{lstlisting}
    \end{exampleblock}
\end{frame}

\section{Reduccion: Suma y Media por Ejes}

\begin{frame}{Reduccion de tensores}
    \begin{itemize}
        \item Consiste en aplicar una operacion (suma, media, maximo) a lo largo de uno o varios ejes.
        \item El resultado elimina esos ejes (o los mantiene con tamaño 1 si usamos \texttt{keepdims}).
    \end{itemize}
    \begin{block}{Suma total}
        \[ \sum_{i,j} X_{ij} \]
    \end{block}
    \begin{block}{Suma por filas (axis=1)}
        \[ (\sum_{j} X_{ij})_i \]
    \end{block}
\end{frame}

\begin{frame}[fragile]{Ejemplos con NumPy}
\begin{lstlisting}
M = np.array([[1,2,3],[4,5,6]])
print(M.sum())
print(M.sum(axis=0))
print(M.sum(axis=1))
print(M.mean(axis=0))
print(M.max(axis=1))
\end{lstlisting}
\end{frame}

\begin{frame}{Aplicacion: Perdida promedio en un batch}
    Durante el entrenamiento, se calcula la perdida para cada muestra y se promedia:
    \[ \text{loss} = \frac{1}{B}\sum_{i=1}^{B} \ell(y_i, \hat{y}_i) \]
    En NumPy: \texttt{loss = np.mean(losses)}.
\end{frame}

\section{Producto Punto}

\begin{frame}{Definicion de producto punto}
    Para vectores $\mathbf{u}, \mathbf{v} \in \mathbb{R}^n$:
    \[ \mathbf{u}\cdot\mathbf{v} = \sum_{i=1}^n u_i v_i = \|\mathbf{u}\| \|\mathbf{v}\| \cos\theta \]
    \begin{itemize}
        \item Mide la proyeccion de un vector sobre otro.
        \item Geometricamente, $\theta$ es el angulo entre ellos.
    \end{itemize}
\end{frame}

\begin{frame}{Producto matriz-vector y matriz-matriz}
    \begin{block}{Matriz-vector}
        \[ (\mathbf{A}\mathbf{x})_i = \sum_{j=1}^{n} A_{ij} x_j \]
    \end{block}
    \begin{block}{Matriz-matriz}
        \[ (\mathbf{A}\mathbf{B})_{ij} = \sum_{k=1}^{m} A_{ik} B_{kj} \]
    \end{block}
\end{frame}

\begin{frame}[fragile]{Ejemplos con NumPy}
\begin{lstlisting}
u = np.array([1, -2, 3])
v = np.array([4, 0, -1])
print(np.dot(u, v))

A = np.array([[1,2],[3,4]])
x = np.array([2,1])
print(A @ x)

B = np.array([[0,1],[1,0]])
print(A @ B)
print(B @ A)
\end{lstlisting}
\end{frame}

\begin{frame}{Similitud de coseno}
    \[ \cos(\mathbf{u},\mathbf{v}) = \frac{\mathbf{u}\cdot\mathbf{v}}{\|\mathbf{u}\|\,\|\mathbf{v}\|} \]
    \begin{exampleblock}{Ejemplo}
        $\mathbf{u}=(1,0,-1)$, $\mathbf{v}=(2,1,0)$:
        \[ \mathbf{u}\cdot\mathbf{v}=2,\quad \|\mathbf{u}\|=\sqrt{2},\quad \|\mathbf{v}\|=\sqrt{5},\quad \cos\theta=\frac{2}{\sqrt{10}}\approx0.632 \]
    \end{exampleblock}
\end{frame}

\begin{frame}{Aplicacion en Transformers: Atencion}
    \begin{block}{Dot-Product Attention}
        \[ \text{Attention}(Q,K,V) = \text{softmax}\left(\frac{QK^\top}{\sqrt{d_k}}\right)V \]
    \end{block}
    \begin{itemize}
        \item $Q,K,V$ son matrices.
        \item $QK^\top$ calcula la similitud entre consultas y claves.
    \end{itemize}
\end{frame}

\section{Ejemplo Integrado}

\begin{frame}[fragile]{Implementacion paso a paso}
\begin{lstlisting}
import numpy as np
batch = np.random.rand(2,3,5)

media = batch.mean(axis=2, keepdims=True)
std = batch.std(axis=2, keepdims=True)
batch_norm = (batch - media) / (std + 1e-8)

def cos_sim_matrix(seq):
    norms = np.linalg.norm(seq, axis=1, keepdims=True)
    sim = np.dot(seq, seq.T) / (norms @ norms.T + 1e-8)
    return sim

sim_matrices = np.array([cos_sim_matrix(seq) for seq in batch_norm])
print(sim_matrices.shape)
\end{lstlisting}
\end{frame}

\begin{frame}
    \centering
    \Huge \textbf{Gracias!}

\end{frame}

\end{document}